\chapter*{Висновки}
\addcontentsline{toc}{chapter}{Висновки}

% ============================================================
%  ВИСНОВКИ -- узагальнення 8 наукових результатів (1:1 з новизною),
%  їх практичне значення, обмеження та напрями подальших досліджень.
%  Кожен пункт = один отриманий результат + чим саме підтверджено.
% ============================================================

У дисертації розв'язано науково-прикладну проблему приватного контекстного захисту учасників цифрової комунікації від ризикових інформаційних впливів в умовах наскрізного шифрування та постквантового переходу. Одержано такі основні результати.

\begin{enumerate}
  \item Запропоновано AURA Runtime-орієнтовану архітектуру приватного контекстного захисту, у якій оцінювання ризику, policy-aware рішення, постквантово-посилена автентифікація, гібридне наскрізне шифрування та серверно-сліпий канал нагляду є рівнями однієї інформаційної технології без перенесення відкритого тексту на сервер. \emph{[Уточнити після розділу 2.]}
  \item Розроблено метод контекстної оцінки інформаційного ризику, що враховує контакт, історію розмови, часову динаміку, повторюваність та латентні стани маніпуляції і формує структурований risk signal. \emph{[Розділ 3.]}
  \item Удосконалено метод policy-aware прийняття safety-рішень з опрацюванням невизначеності, профілем захисту учасника та privacy-safe аудитом. \emph{[Розділ 4.]}
  \item Удосконалено метод постквантово-посиленої OPAQUE-автентифікації з гібридним DH/KEM-узгодженням сесії та операційним розмежуванням \texttt{oprf\_seed}, що підтверджено формально та реалізаційно. \emph{[Розділ 5.]}
  \item Удосконалено метод гібридного постквантового наскрізного шифрування з KEM-підсиленим кроком «храповика» та scoped-формальними доказами. \emph{[Розділ 6.]}
  \item Запропоновано метод серверно-сліпого каналу нагляду, що передає safety-сигнали guardian endpoint-у без розкриття відкритого тексту серверу. \emph{[Розділ 7.]}
  \item Удосконалено модель меж компрометації та методологію приватного аудиту й release-gated оцінювання зі статусами \texttt{PASS}, \texttt{FAIL}, \texttt{INSUFFICIENT\_SUPPORT}, \texttt{BLOCKED}. \emph{[Розділи 2, 8.]}
  \item Розроблено інформаційну технологію AURA/Ecliptix, яка реалізує запропоновані методи з відтворюваною перевіркою. \emph{[Розділ 8.]}
\end{enumerate}

\noindent\textbf{Практичне значення.} \emph{[Узагальнити придатність результатів та впровадження.]}

\medskip
\noindent\textbf{Обмеження та напрями подальших досліджень.} \emph{[Чесно зазначити межі: відсутність клінічної валідації, класичний характер підписів, scope формальних доказів; окреслити подальші напрями.]}
